% Shared preamble for one-page research-statement cover letters
% Build: pdflatex <letter>.tex
\documentclass[10pt,letterpaper]{article}
\usepackage[T1]{fontenc}
\usepackage[utf8]{inputenc}
\usepackage[sc]{mathpazo}          % Palatino, real small caps
\linespread{1.02}
\usepackage[left=0.9in,right=0.9in,top=0.6in,bottom=0.55in]{geometry}
\usepackage{microtype}
\usepackage{xcolor}
\definecolor{accent}{RGB}{22,48,92}    % deep navy
\definecolor{muted}{gray}{0.35}
\usepackage[hidelinks]{hyperref}
\urlstyle{same}

\pagestyle{empty}
\setlength{\parindent}{0pt}
\setlength{\parskip}{0.4\baselineskip}

% Inline venue citation, research-statement style: \venue{EuroSys '16}
\newcommand{\venue}[1]{{\small[#1]}}

% Reference list keyed by the SAME tag the inline \venue{} prints, so a reader
% matches a citation to its entry without a separate numbering scheme.
\newcommand{\refhead}[1]{%
  {\large\bfseries\color{accent}#1}\\[-5pt]
  {\color{accent}\rule{\textwidth}{0.6pt}}\par\vspace{1pt}%
}
\newenvironment{references}{\small\setlength{\parskip}{3pt}}{\par}
% Tags vary in width (``CCS '23'' vs ``USENIX ;login: '16''), so the key runs
% inline with a hanging indent rather than in a fixed-width box that would overflow.
\newcommand{\refitem}[2]{%
  \noindent\hangindent=1.8em\hangafter=1
  \textbf{[#1]}\hspace{0.4em}\ignorespaces #2\par%
}

% Small-caps run-in paragraph lead: \lead{The agenda.}
\newcommand{\lead}[1]{\textsc{\color{accent}#1}\hspace{0.5em}}

% Letterhead: name + affiliations left, contact right, accent rule
\newcommand{\letterhead}{%
  \noindent
  \begin{minipage}[b]{0.63\textwidth}
    {\LARGE\bfseries Vaggelis Atlidakis}\\[3pt]
    {\small\color{muted}Assistant Professor, Informatics \& Telecommunications,
      University of Athens}%
  \end{minipage}\hfill
  \begin{minipage}[b]{0.33\textwidth}
    \raggedleft\small
    \href{mailto:v.atlidakis@gmail.com}{v.atlidakis@gmail.com}\\
    \href{https://vatlidak.github.io/}{vatlidak.github.io}\\
    +30\,698\,992\,9982%
  \end{minipage}\\[7pt]
  {\color{accent}\rule{\textwidth}{1pt}}%
}

% Addressee (left) and date (right): \dateline{Trail of Bits}{Re: ...}{August 7, 2026}
\newcommand{\dateline}[3]{%
  \noindent
  \begin{minipage}[t]{0.63\textwidth}
    \textbf{#1}\\{\small\color{muted}#2}%
  \end{minipage}\hfill
  \begin{minipage}[t]{0.33\textwidth}
    \raggedleft #3%
  \end{minipage}\par
}

% Centered statement title + subtitle
\newcommand{\lettertitle}[2]{%
  \begin{center}
    {\large\bfseries\color{accent} #1}\\[3pt]
    {\small\itshape #2}
  \end{center}\vspace{-2pt}%
}

% Signature block
\newcommand{\signature}{%
  \smallskip
  Sincerely,\\[3pt]
  {\bfseries Vaggelis Atlidakis},\\[1pt]
  {\small\color{muted}%
  Assistant Professor, Informatics \& Telecommunications\\
  National and Kapodistrian University of Athens\\
  M.Sc.\ \textperiodcentered\ M.Phil.\ \textperiodcentered\ Ph.D.\ in Computer Science, Columbia University\\
  Formerly: Brown University \textperiodcentered\ Microsoft Research \textperiodcentered\ CERN}\par
}

\hypersetup{
  pdftitle={Vaggelis Atlidakis -- Research Statement / Cover Letter -- Trail of Bits},
  pdfauthor={Vaggelis Atlidakis},
}
\begin{document}

\letterhead
\vspace{2pt}
\dateline{Trail of Bits}{Re: General application}{August 7, 2026}

\lettertitle{Systems Abstractions for Agentic Workloads}{Research statement accompanying a general application}

Dear Trail of Bits team,

Ten years ago, my collaborators and I measured how modern applications actually
use POSIX abstractions and found them aging out of relevance: applications had
outgrown them, new abstractions were accreting as ad-hoc user-space libraries, and
applications leaned heavily on extension interfaces such as \texttt{ioctl}, which
the kernel cannot meaningfully inspect---breeding mismatches, inefficiencies, and
security risks \venue{EuroSys~'16}. We put it bluntly at the time: ``POSIX has
become outdated'' \venue{USENIX {\itshape;login:}~'16}. The same story is repeating
today, faster, and over interfaces that reach the filesystem, the network, and the
shell directly. The new workload class is
agentic---LLM-driven agents that plan, spawn sub-agents, invoke tools, and
execute freshly generated code---and it runs on abstractions designed for human
users and long-lived services. Every agent runtime re-implements identity,
sandboxing, resource control, and audit in user space, each one differently:
the new libraries-over-POSIX. My agenda is to establish, by measurement,
precisely what fails in traditional systems abstractions under agentic
workloads---and to design the new ones that make these workloads portable,
secure, and efficient. Those abstractions---the POSIX of the agentic era---do
not exist yet; I intend to help define them.

\lead{The agenda.} I plan to rerun a playbook I know works.
(i)~\emph{Measure}: a workload study of agent runtimes and tool protocols
against today's OS abstractions---the old, the new, and the missing, redone
for agents.
(ii)~\emph{Design}: the primitives the data points to---task-scoped identity
rather than user-scoped; least privilege refined dynamically as agents switch
tasks (SysXCHG, which re-instantiates system-call filters on \texttt{execve},
is a direct ancestor \venue{CCS~'23}); isolation cheap enough to wrap every
tool call; provenance as a first-class object rather than a log line.
(iii)~\emph{Test}: agent interfaces need what REST APIs got---automated,
stateful, feedback-driven testing---and I built that the first time as RESTler
\venue{ICSE~'19}.

\lead{The record.} I have worked both sides of this agenda: discovering
vulnerabilities automatically, and building the mechanisms that contain them.
RESTler, the first stateful REST API fuzzer, found hundreds of bugs in
production Azure, O365, and Bing services, open-sourced and patented
\venue{ICSE~'19, ICST~'20}; IvySyn uncovered 61 previously-unknown
vulnerabilities (39 CVEs) in TensorFlow and PyTorch
\venue{USENIX Sec~'23}. On defense, FineIBT---fine-grain, hardware-assisted
control-flow integrity---was adopted by Intel and upstreamed, in a variant, to
the Linux kernel \venue{RAID~'23}, and Eclipse blocks speculative abuse of
memory errors \venue{CCS~'24}. I was an NSF/CRA Computing Innovation Fellow,
and I ship production software in industry alongside my tenure-track post in
Athens; I serve on the 2026 CCS, EuroSys, and RAID program committees.

\lead{Why Trail of Bits.} Buttercup's second-place finish at DARPA's AIxCC
settled whether agents can find and patch real vulnerabilities
\venue{AIxCC~'25}. What remains unsettled is the substrate they run on: every
cyber-reasoning system still assembles its own orchestration, sandboxing, and
budget enforcement. That missing layer is a research program with tool-shaped
output, and I want to pursue it the way Trail of Bits does---in the open, as
artifacts people can run, which is how nearly everything I build already ships.
This letter answers no specific posting: it proposes a direction of work, and I
would like to find out where in your research or assurance practice it belongs.

I would welcome a conversation.

\signature

\newpage

{\small\color{muted}Vaggelis Atlidakis \,\textperiodcentered\, Systems Abstractions
for Agentic Workloads}\\[7pt]
\refhead{References}

\begin{references}

\refitem{EuroSys~'16}{V.~Atlidakis, J.~Andrus, R.~Geambasu, D.~Mitropoulos, and
J.~Nieh. POSIX Abstractions in Modern Operating Systems: The Old, the New, and
the Missing. \textit{EuroSys}, 2016, pp.~19:1--19:17.
\href{https://doi.org/10.1145/2901318.2901350}{doi:10.1145/2901318.2901350}.
Artifact: \href{https://github.com/columbia/libtrack}{github.com/columbia/libtrack}}

\refitem{USENIX {\itshape;login:}~'16}{V.~Atlidakis, J.~Andrus, R.~Geambasu,
D.~Mitropoulos, and J.~Nieh. POSIX Has Become Outdated. \textit{USENIX ;login:},
41(3), Fall 2016.
\href{https://www.usenix.org/publications/login/fall2016/atlidakis}{usenix.org/publications/login/fall2016/atlidakis}}

\refitem{CCS~'23}{A.~J.~Gaidis, V.~Atlidakis, and V.~P.~Kemerlis. SysXCHG:
Refining Privilege with Adaptive System Call Filters. \textit{ACM CCS}, 2023,
pp.~1964--1978.
Artifact: \href{https://gitlab.com/brown-ssl/sysxchg}{gitlab.com/brown-ssl/sysxchg}}

\refitem{ICSE~'19}{V.~Atlidakis, P.~Godefroid, and M.~Polishchuk. RESTler:
Stateful REST API Fuzzing. \textit{ICSE}, 2019, pp.~748--758. US Patent
10,977,161. Open source:
\href{https://github.com/microsoft/restler-fuzzer}{github.com/microsoft/restler-fuzzer}}

\refitem{ICST~'20}{V.~Atlidakis, P.~Godefroid, and M.~Polishchuk. Checking
Security Properties of Cloud Service REST APIs. \textit{IEEE ICST}, 2020,
pp.~387--397.}

\refitem{USENIX Sec~'23}{N.~Christou, D.~Jin, V.~Atlidakis, B.~Ray, and
V.~P.~Kemerlis. IvySyn: Automated Vulnerability Discovery in Deep Learning
Frameworks. \textit{USENIX Security}, 2023, pp.~2383--2400.
Artifact: \href{https://gitlab.com/brown-ssl/ivysyn}{gitlab.com/brown-ssl/ivysyn}}

\refitem{RAID~'23}{A.~J.~Gaidis, J.~Moreira, K.~Sun, A.~Milburn, V.~Atlidakis,
and V.~P.~Kemerlis. FineIBT: Fine-grain Control-flow Enforcement with Indirect
Branch Tracking. \textit{RAID}, 2023, pp.~527--546.
Artifact: \href{https://gitlab.com/brown-ssl/fineibt}{gitlab.com/brown-ssl/fineibt}}

\refitem{CCS~'24}{N.~Christou, A.~J.~Gaidis, V.~Atlidakis, and V.~P.~Kemerlis.
Eclipse: Preventing Speculative Memory-error Abuse with Artificial Data
Dependencies. \textit{ACM CCS}, 2024, pp.~3913--3927.
Artifact: \href{https://gitlab.com/brown-ssl/eclipse}{gitlab.com/brown-ssl/eclipse}}

\refitem{AIxCC~'25}{Trail of Bits. Buttercup: an open-source AI Cyber Reasoning
System. Second place, DARPA AI Cyber Challenge finals, August 2025.
\href{https://github.com/trailofbits/buttercup}{github.com/trailofbits/buttercup}}

\end{references}

\end{document}
